% !TeX root = ../main.tex

\xchapter{基于自回归的人体姿态重构与预测}{Algorithm and Code}

在实现了对人体宏观活动的高精度识别后，本章将研究视角进一步延伸至微观的人体动作细节，深入探讨基于惯性测量单元（IMU）的人体姿态重构与预测问题。第4.1节首先对人体姿态重构与预测任务进行了详细描述，回顾了现有非光学捕捉技术在实时性与预测能力上的不足，明确了引入自回归机制的研究动机。第4.2节介绍了适用于该任务的数据集构建细节，确保了数据的多样性与质量。第4.3节系统地阐述了本文提出的基于自回归机制的人体姿态重构与预测方法，详细说明了模型中IMU特征提取器、姿态重构器与IMU预测器的协同工作原理，以及如何通过自回归闭环实现对当前姿态的重构与未来动作序列的精准预测。第4.4节展示了广泛的实验结果与深入分析，除了基础的评价指标与实验环境介绍外，还重点涵盖了数据实验、骨干网络选择、标准化消融实验、跨域适应性分析以及最大预测长度对模型性能影响的探讨。通过这些实验，全面评估并验证了模型在实现厘米级精度重构与长时间序列预测方面的优越性能。

\xsection{问题描述}{Algorithm}

在数字化浪潮的推动下，元宇宙与数字孪生技术正逐步从概念走向落地，而人体姿态重构作为连接物理世界人类本体与虚拟数字空间化身的核心纽带，其技术成熟度直接决定了人机交互的沉浸感与真实感。相比于受限于视场角、光照条件及隐私敏感性的计算机视觉方案，基于稀疏惯性测量单元（IMU）的姿态重构技术凭借其全天候可用、不受遮挡影响、佩戴轻便及低算力成本等原生优势，已成为虚拟现实交互、远程医疗康复、云端体感游戏以及工业人机协作等领域的首选技术路径。然而，尽管近年来基于深度学习的惯性姿态重构方法在标准数据集上的重建精度已取得显著突破，但在面向高保真、低延迟的实时交互应用时，现有的主流技术方案仍面临着难以逾越的理论瓶颈与工程挑战。这些挑战并非源于传感器精度的不足，而是源于传统建模范式对人体运动规律的认知局限，主要集中在时序预测能力的缺失、运动学一致性的破坏以及对抗干扰鲁棒性的不足三个深层次维度。

第一，现有主流的判别式回归模型本质上是一种静态映射机制，缺乏对未来动作趋势的主动预测能力，无法从根本上解决实时交互系统中的端到端延迟痛点。在当前的虚拟现实云渲染、远程手术机器人控制等典型应用场景中，数据从传感器采集、无线网络传输、云端解算到终端画面渲染，整个链路不可避免地会引入数十至数百毫秒的系统延迟。现有的技术路线大多采用判别式建模范式，即训练一个深度神经网络来拟合从当前时刻或历史窗口的IMU观测值到当前时刻人体姿态参数的映射函数。这种所见即所得的瞬时重构机制虽然在离线评估中能够获得较高的逐帧精度，但在在线实时应用中却存在致命缺陷：终端渲染出的虚拟人动作永远滞后于用户的真实物理动作。根据控制理论与认知心理学研究，这种视觉反馈与前庭系统感知之间的时序错位（即视觉-前庭冲突）不仅会严重破坏用户的沉浸感与临场感，在高速运动或剧烈转头等高动态场景下，甚至会引发严重的晕动症与操作失误。因此，迫切需要一种能够打破时间壁垒，从被动重构历史转向主动预测未来的新型建模范式，通过精准预测未来的姿态序列来从算法层面抵消系统物理延迟，实现零延迟甚至负延迟的极致交互体验。

第二，传统的帧级独立重构或短时序回归方法忽略了人体运动的内在连续性与生物力学约束，难以保证重构姿态序列的长时序平滑性与物理合理性。人体运动是一个受骨骼长度、关节活动度及肌肉动力学严格约束的连续物理过程，具有高度的时间相关性与状态依赖性。然而，现有的许多方法为了追求推理速度，往往将姿态重构简化为帧级独立的回归问题，或者仅利用极短的历史窗口进行特征提取，从而切断了姿态序列之间的长程依赖关系。这种做法导致模型对微机电系统传感器固有的高频噪声与零偏漂移极其敏感，重构出的动作序列往往伴随着非自然的抖动、高频振荡或违反关节角度生理限制的突变。虽然工业界常采用卡尔曼滤波或巴特沃斯低通滤波器进行后处理平滑，但这本质上是一种拆东墙补西墙的妥协策略，滤波器的引入会抹杀动作的高频细节（如快速出拳或瞬时急停），并不可避免地引入额外的相位延迟，进一步恶化了系统的实时响应能力。如何内嵌地建模人体运动的动力学演变规律，在生成过程中自动约束姿态的时空一致性，是提升重构质量的关键所在。

第三，针对传感器数据丢失、信号漂移或异常输入的鲁棒性生成能力不足，限制了现有系统在复杂电磁环境下的可靠性。在实际的穿戴应用环境中，商用IMU设备常面临蓝牙信号丢包、周围铁磁物质干扰导致的磁力计失准以及长时间运行后的零偏漂移等问题。传统的回归模型通常假设输入数据是完整且服从训练集分布的，缺乏对数据分布本身的生成式建模能力。一旦输入数据出现缺失、跳变或分布外噪声，回归网络往往会输出不可预测的崩塌结果，导致虚拟人肢体扭曲或瞬移。相比之下，生成式模型能够学习人体运动流形的潜在分布，天然具备序列补全与去噪的能力，能够利用历史记忆维持动作的连贯性，具有更高的容错上限。

鉴于上述深层次的理论与应用挑战，本研究认为，人体姿态重构不应仅仅被视为一个简单的信号到信号的映射问题，而应被重构为一个更为宏大的序列生成与预测问题。基于此，本文提出了基于自回归机制的人体姿态重构与预测方法。该方法摒弃了传统的直接映射思路，创新性地引入了IMU预测器与闭环反馈机制，构建了一个能够自我演进的生成式架构。通过联合建模当前姿态与未来传感器数据的条件概率分布，该模型不仅能够利用历史信息平滑当前的姿态估计（重构任务），还能根据学习到的运动趋势推演未来的IMU数据并进而生成未来的姿态（预测任务）。这种将重构与预测统一在同一自回归框架下的设计，不仅显著提升了系统在复杂环境下的鲁棒性与平滑性，更为构建下一代低延迟、高保真的元宇宙交互系统提供了全新的技术路径与理论支撑。

考虑到实际使用需求，即长算法分段和注释功能，模板不再使用 \clist[columbiablue]{algorithm2e} 作为默认算法排版工具，改为调用 \clist[columbiablue]{algpseudocode}。注意，这两个宏包相互冲突。算法的具体编写方式请自行查阅宏包，即\clist{texdoc algpseudocode}。

\begin{texcode}[]{}
\begin{algorithm}[H]
  \caption{How to write algorithms}
  \begin{algorithmic}[1]
      \State \textbf{Input:} this text
      \State \textbf{Output:} how to write algorithm with \LaTeX2e \Comment{Just a comment}
      \State initialization
      \While{not at end of this document}
        \State read current
        \If{understand}
          \State go to next section
          \State current section becomes this one
        \Else{$k\geq k^\text{max}$}
          \State go back to the beginning of current section
        \EndIf
      \EndWhile
  \end{algorithmic}
\end{algorithm}
\end{texcode}

\begin{texcode}[]{}
  \begin{algorithm}[H]
    \caption{Simulation-optimization heuristic}\label{algorithm:SOH}
    \begin{algorithmic}[1]
        \State \textbf{Data:} current period $t$, initial inventory $I_{t-1}$, initial capital $B_{t-1}$, demand samples
        \State \textbf{Output:} Optimal order quantity $Q^{\ast}_{t}$
        \State $r\leftarrow t$
        \State $\Delta B^{\ast}\leftarrow -\infty$
        \While{$\Delta B\leq \Delta B^{\ast}$ and $r\leq T$}
          \State $Q\leftarrow\arg\max_{Q\geq 0}\Delta B^{Q}_{t,r}(I_{t-1},B_{t-1})$
          \State $\Delta B\leftarrow \Delta B^{Q}_{t,r}(I_{t-1},B_{t-1})/(r-t+1)$
          \If{$\Delta B\geq \Delta B^{\ast}$}
            \State $Q^{\ast}\leftarrow Q$
            \State $\Delta B^{\ast}\leftarrow \Delta B$
          \EndIf
          \State $r\leftarrow r+1$
        \EndWhile
    \end{algorithmic}
  \end{algorithm}
\end{texcode}

\xsection{数据集构建}{Input codes}

从现有文件导入源码，使用 \clist[columbiablue]{lstlisting} 宏包，本模板预先定义了一基本格式 \clist{sty_basic} 以供使用，可以自定义新的样式。

\begin{texcodeonly}
\lstinputlisting[style=sty_basic,
  language=matlab,
  caption={标准粒子群算法},
  label={listing:a_label}]{./Codes/pso.m}
\end{texcodeonly}
\lstinputlisting[style=sty_basic,
                    language=matlab,
                    caption={标准粒子群算法},
                    label={listing:a_label}]{./Codes/pso.m}


\xsection{方法描述}{Input codes}

\begin{figure}[h]
  \centering
  \includegraphics[width=\textwidth]{姿态重构模型架构图.jpg}
  \caption{人体姿态重构与预测模型架构图}
\end{figure}

\xsection{实验与结果分析}{Input codes}

\xsubsection{评价指标}{Level of this document}

\xsubsection{实验环境}{Level of this document}

\xsubsection{数据实验结果分析}{Level of this document}

\subsubsection{数据分割实验}{Ref equations}

\begin{table}[H]
  \centering
  \begin{tblr}{
    colspec={X[c] X[c] X[c]},
  }
  \toprule
  IMU输入长度 & reconstruction mpjpe & prediction mpjpe \\
  \midrule

  15  & 78.1419 & 81.7890 \\
  30  & 63.0527 & 69.1466 \\
  45  & 56.0555 & 63.0849 \\
  60  & 50.7079 & 56.6010 \\
  75  & 48.6763 & 54.1706 \\
  90  & 48.7303 & 53.8068 \\
  105 & 47.3098 & 51.6719 \\
  120 & 47.4721 & 52.1078 \\
  135 & 48.6094 & 53.8773 \\
  150 & 47.9178 & 52.9400 \\

  \bottomrule
  \end{tblr}
\end{table}

\subsubsection{设备消融实验}{Ref equations}

\begin{table}[H]
  \centering
  \begin{tblr}{
    colspec={X[c] X[c] X[c]},
  }
  \toprule
  消除设备 & reconstruction mpjpe & prediction mpjpe \\
  \midrule

  None      & 50.7079 & 56.6010 \\
  0         & 52.2946 & 58.6593 \\
  1         & 51.2945 & 57.8350 \\
  2         & 50.5097 & 56.7976 \\
  3         & 51.3100 & 57.8461 \\
  4         & 50.6537 & 56.2470 \\
  2 and 4   & 52.5957 & 58.7247 \\

  \bottomrule
  \end{tblr}
\end{table}

\xsubsection{骨干网络选择实验结果分析}{Level of this document}

\begin{table}[H]
  \centering
  \begin{tblr}{
    colspec={X[c] X[c] X[c] X[c]},
    % row{1} = {font=\bfseries}
  }
  \toprule
  主干网络1 & 主干网络2 & reconstruction mpjpe & prediction mpjpe \\
  \midrule
  
  resnet & transformer & 47.8203 & 54.1138 \\
  resnet & lstm        & 49.4344 & 55.9392 \\
  resnet & gru         & 48.9044 & 54.9700 \\
  resnet & mamba       & 76.8240 & 81.0473 \\
  
  \bottomrule
  \end{tblr}
\end{table}

\xsubsection{标准化消融实验结果分析}{Level of this document}

\begin{table}[H]
  \centering
  \begin{tblr}{
    colspec={X[c] X[c] X[c] X[c]},
    % row{1} = {font=\bfseries}
  }
  \toprule
  IMU标准化 & Pose标准化 & reconstruction mpjpe & prediction mpjpe \\
  \midrule
  
  否 & 否 & 60.7588 & 74.1964 \\
  
  否 & 是 & 79.7920 & 86.5976 \\
  
  是 & 否 & 58.5465 & 70.5447 \\
  
  是 & 是 & 50.6833 & 56.4806 \\
  
  \bottomrule
  \end{tblr}
\end{table}

\xsubsection{跨域实验结果分析}{Level of this document}

\subsubsection{跨场景实验}{Ref equations}

\begin{table}[H]
  \centering
  \begin{tblr}{
    colspec={X[c] X[c] X[c]},
    % row{1} = {font=\bfseries}
  }
  \toprule
  测试场景 & reconstruction mpjpe & prediction mpjpe \\
  \midrule

  livingroom & 204.4093 & 226.1702 \\
  office     & 147.1609 & 155.1837 \\
  kitchen    & 138.3596 & 143.9825 \\

  \bottomrule
  \end{tblr}
\end{table}

\subsubsection{跨人实验}{Ref equations}

\begin{table}[H]
  \centering
  \begin{tblr}{
    colspec={X[c] X[c] X[c]},
  }
  \toprule
  测试志愿者 & reconstruction mpjpe & prediction mpjpe \\
  \midrule

  0  & 93.9199  & 100.6215 \\
  1  & 117.3705 & 123.4515 \\
  2  & 86.6879  & 91.1092 \\
  3  & 102.1780 & 108.2353 \\
  4  & 126.7194 & 138.4886 \\
  5  & 104.7901 & 113.8852 \\
  6  & 100.7176 & 106.4436 \\
  7  & 108.6058 & 118.8354 \\
  8  & 114.5809 & 122.3770 \\
  9  & 98.3641  & 106.3886 \\
  10 & 97.3505  & 104.0960 \\
  11 & 105.0606 & 112.9244 \\
  12 & 85.6985  & 89.5146 \\
  13 & 95.7136  & 102.6839 \\
  14 & 103.5153 & 110.9676 \\
  15 & 115.2117 & 125.8084 \\

  \bottomrule
  \end{tblr}
\end{table}

\xsubsection{最大预测长度实验结果分析}{Level of this document}

\begin{table}[H]
  \centering
  \begin{tblr}{
    colspec={X[c] X[c]},
  }
  \toprule
  预测长度区间 & mpjpe \\
  \midrule

  {[-15, 0)}   & \texttt{56.1941} \\
  {[0, 15)}    & \texttt{60.9056} \\
  {[15, 30)}   & \texttt{65.6229} \\
  {[30, 45)}   & \texttt{77.7804} \\
  {[45, 60)}   & \texttt{88.2202} \\
  {[60, 75)}   & \texttt{94.8827} \\
  {[75, 90)}   & \texttt{101.9103} \\
  {[90, 105)}  & \texttt{108.3201} \\
  {[105, 120)} & \texttt{113.4193} \\
  {[120, 135)} & \texttt{119.1812} \\
  {[135, 150)} & \texttt{125.6723} \\

  \bottomrule
  \end{tblr}
\end{table}